\subsection{Euler's pentagonal number theorem}

The following definition looks somewhat quaint; why define a notation for a
specific quadratic function?

\begin{definition}
\label{def.pars.pent-num}For any $k\in\mathbb{Z}$, define a nonnegative
integer $w_{k}\in\mathbb{N}$ by%
\[
w_{k}=\dfrac{\left(  3k-1\right)  k}{2}.
\]
This is called the $k$\emph{-th pentagonal number}.
\end{definition}

Here is a table of these pentagonal numbers:%
\[%
\begin{tabular}
[c]{|c||c|c|c|c|c|c|c|c|c|c|c|c|c|}\hline
$k$ & $\cdots$ & $-5$ & $-4$ & $-3$ & $-2$ & $-1$ & $0$ & $1$ & $2$ & $3$ &
$4$ & $5$ & $\cdots$\\\hline
$w_{k}$ & $\cdots$ & $40$ & $26$ & $15$ & $7$ & $2$ & $0$ & $1$ & $5$ & $12$ &
$22$ & $35$ & $\cdots$\\\hline
\end{tabular}
\ \ \ \ .
\]
Note that $w_{k}$ really is a nonnegative integer for any $k\in\mathbb{Z}$
(check this!). The name \textquotedblleft pentagonal numbers\textquotedblright%
\ is historically motivated (see
\href{https://en.wikipedia.org/wiki/Pentagonal_number}{the Wikipedia page} for
details); the only thing we need to know about them (beside their definition)
is the fact that they are nonnegative integers and grow quadratically with $n$
in both directions (i.e., when $n\rightarrow\infty$ and when $n\rightarrow
-\infty$). The latter fact ensures that the infinite sum $\sum_{k\in
\mathbb{Z}}\left(  -1\right)  ^{k}x^{w_{k}}$ is a well-defined FPS in
$\mathbb{Z}\left[  \left[  x\right]  \right]  $. Rather surprisingly, this
infinite sum coincides with a particularly simple infinite product:

\begin{theorem}
[Euler's pentagonal number theorem]\label{thm.pars.pent}We have%
\[
\prod_{k=1}^{\infty}\left(  1-x^{k}\right)  =\sum_{k\in\mathbb{Z}}\left(
-1\right)  ^{k}x^{w_{k}}.
\]

\end{theorem}

Let us write this out concretely:%
\begin{align*}
&  \prod_{k=1}^{\infty}\left(  1-x^{k}\right) \\
&  =\sum_{k\in\mathbb{Z}}\left(  -1\right)  ^{k}x^{w_{k}}\\
&  =\cdots+x^{w_{-4}}-x^{w_{-3}}+x^{w_{-2}}-x^{w_{-1}}+x^{w_{0}}-x^{w_{1}%
}+x^{w_{2}}-x^{w_{3}}+x^{w_{4}}-x^{w_{5}}\pm\cdots\\
&  =\cdots+x^{26}-x^{15}+x^{7}-x^{2}+1-x+x^{5}-x^{12}+x^{22}-x^{35}\pm\cdots\\
&  =1-x-x^{2}+x^{5}+x^{7}-x^{12}-x^{15}+x^{22}+x^{26}\pm\cdots.
\end{align*}


We will prove Theorem \ref{thm.pars.pent} in the next section (as a particular
case of Jacobi's Triple Product Identity).\footnote{See \cite{Bell06} for the
history of Theorem \ref{thm.pars.pent}.} First, let us use it to derive the
following recursive formula for the partition numbers $p\left(  n\right)  $:

\begin{corollary}
\label{cor.pars.pn-rec}For each positive integer $n$, we have%
\begin{align*}
p\left(  n\right)   &  =\sum_{\substack{k\in\mathbb{Z};\\k\neq0}}\left(
-1\right)  ^{k-1}p\left(  n-w_{k}\right) \\
&  =p\left(  n-1\right)  +p\left(  n-2\right)  -p\left(  n-5\right)  -p\left(
n-7\right) \\
&  \ \ \ \ \ \ \ \ \ \ +p\left(  n-12\right)  +p\left(  n-15\right)  -p\left(
n-22\right)  -p\left(  n-26\right)  \pm\cdots.
\end{align*}

\end{corollary}

\begin{proof}
[Proof of Corollary \ref{cor.pars.pn-rec} using Theorem \ref{thm.pars.pent}%
.]We have%
\[
\sum_{m\in\mathbb{N}}p\left(  m\right)  x^{m}=\sum_{n\in\mathbb{N}}p\left(
n\right)  x^{n}=\prod\limits_{k=1}^{\infty}\dfrac{1}{1-x^{k}}%
\]
(by Theorem \ref{thm.pars.main-gf}) and%
\[
\sum_{k\in\mathbb{Z}}\left(  -1\right)  ^{k}x^{w_{k}}=\prod\limits_{k=1}%
^{\infty}\left(  1-x^{k}\right)
\]
(by Theorem \ref{thm.pars.pent}). Multiplying these two equalities, we obtain%
\begin{align}
\left(  \sum_{m\in\mathbb{N}}p\left(  m\right)  x^{m}\right)  \cdot\left(
\sum_{k\in\mathbb{Z}}\left(  -1\right)  ^{k}x^{w_{k}}\right)   &  =\left(
\prod\limits_{k=1}^{\infty}\dfrac{1}{1-x^{k}}\right)  \cdot\left(
\prod\limits_{k=1}^{\infty}\left(  1-x^{k}\right)  \right) \nonumber\\
&  =\prod\limits_{k=1}^{\infty}\underbrace{\left(  \dfrac{1}{1-x^{k}}%
\cdot\left(  1-x^{k}\right)  \right)  }_{=1}\nonumber\\
&  =1. \label{pf.cor.pars.pn-rec.1}%
\end{align}


Now, let us fix a positive integer $n$. We shall compare the $x^{n}%
$-coefficients on both sides of (\ref{cor.pars.pn-rec}).

The $x^{n}$-coefficient on the left hand side of (\ref{cor.pars.pn-rec}) is%
\begin{align*}
\sum_{\substack{m\in\mathbb{N};\\k\in\mathbb{Z};\\m+w_{k}=n}}p\left(
m\right)  \cdot\left(  -1\right)  ^{k}  &  =\sum_{\substack{k\in
\mathbb{Z};\\n-w_{k}\geq0}}p\left(  n-w_{k}\right)  \cdot\left(  -1\right)
^{k}\\
&  \ \ \ \ \ \ \ \ \ \ \ \ \ \ \ \ \ \ \ \ \left(
\begin{array}
[c]{c}%
\text{here, we have replaced }m\text{ by }n-w_{k}\\
\text{in the sum, since the condition }m+w_{k}=n\\
\text{forces }m\text{ to be }n-w_{k}%
\end{array}
\right) \\
&  =\sum_{k\in\mathbb{Z}}p\left(  n-w_{k}\right)  \cdot\left(  -1\right)
^{k}\\
&  \ \ \ \ \ \ \ \ \ \ \ \ \ \ \ \ \ \ \ \ \left(
\begin{array}
[c]{c}%
\text{here, we have extended the range of}\\
\text{summation; this does not change the sum,}\\
\text{since }p\left(  n-w_{k}\right)  =0\text{ whenever }n-w_{k}<0
\end{array}
\right) \\
&  =\sum_{k\in\mathbb{Z}}\left(  -1\right)  ^{k}p\left(  n-w_{k}\right) \\
&  =\underbrace{\left(  -1\right)  ^{0}}_{=1}p\left(  n-\underbrace{w_{0}%
}_{=0}\right)  +\sum_{\substack{k\in\mathbb{Z};\\k\neq0}}\left(  -1\right)
^{k}p\left(  n-w_{k}\right) \\
&  =p\left(  n\right)  +\sum_{\substack{k\in\mathbb{Z};\\k\neq0}}\left(
-1\right)  ^{k}p\left(  n-w_{k}\right)  .
\end{align*}
But the $x^{n}$-coefficient on the right hand side of (\ref{cor.pars.pn-rec})
is $0$ (since $n$ is positive). Hence, comparing the coefficients yields
\[
p\left(  n\right)  +\sum_{\substack{k\in\mathbb{Z};\\k\neq0}}\left(
-1\right)  ^{k}p\left(  n-w_{k}\right)  =0.
\]
Solving this for $p\left(  n\right)  $, we find%
\[
p\left(  n\right)  =-\sum_{\substack{k\in\mathbb{Z};\\k\neq0}}\left(
-1\right)  ^{k}p\left(  n-w_{k}\right)  =\sum_{\substack{k\in\mathbb{Z}%
;\\k\neq0}}\left(  -1\right)  ^{k-1}p\left(  n-w_{k}\right)  .
\]
Corollary \ref{cor.pars.pn-rec} is thus proved.
\end{proof}


\subsection{Jacobi's triple product identity}

\subsubsection{The identity}

Instead of proving Theorem \ref{thm.pars.pent} directly, we shall prove a
stronger result: \emph{Jacobi's triple product identity}. This identity can be
stated as follows:%
\begin{equation}
\prod_{n>0}\left(  \left(  1+q^{2n-1}z\right)  \left(  1+q^{2n-1}%
z^{-1}\right)  \left(  1-q^{2n}\right)  \right)  =\sum_{\ell\in\mathbb{Z}%
}q^{\ell^{2}}z^{\ell}. \label{eq.pars.jtp.jtp}%
\end{equation}
What are $q$ and $z$ here? It appears that (\ref{eq.pars.jtp.jtp}) should be
an identity between multivariate Laurent series (in the indeterminates $q$ and
$z$), but we have never defined such a concept. Multivariate Laurent series
can indeed be defined, but this is not as easy as the univariate case and
involves some choices (see \cite{ApaKau13} for details).

\begin{noncompile}
However, it is not hard to see that any negative power of $z$ in
(\ref{eq.pars.jtp.jtp}) is \textquotedblleft compensated\textquotedblright\ by
a positive power of $q$. For example, the $q^{2n-1}z^{-1}$ term on the left
hand side will always have total degree $\geq0$, and so will the $q^{\ell^{2}%
}z^{\ell}$ term on the right hand side. So this is not as bad as Laurent
series in general. (Commented this out because it is irrelevant and confusing.
The same could be said about $\sum_{\ell\in\mathbb{Z}}q^{-\ell}z^{\ell}$,
which is a bad series.)
\end{noncompile}

A simpler ring in which the identity (\ref{eq.pars.jtp.jtp}) can be placed is
$\left(  \mathbb{Z}\left[  z^{\pm}\right]  \right)  \left[  \left[  q\right]
\right]  $ (that is, the ring of FPSs in the indeterminate $q$ whose
coefficients are Laurent polynomials over $\mathbb{Z}$ in the indeterminate
$z$). In other words, we state the following:

\begin{theorem}
[Jacobi's triple product identity, take 1]\label{thm.pars.jtp1}In the ring
$\left(  \mathbb{Z}\left[  z^{\pm}\right]  \right)  \left[  \left[  q\right]
\right]  $, we have%
\[
\prod_{n>0}\left(  \left(  1+q^{2n-1}z\right)  \left(  1+q^{2n-1}%
z^{-1}\right)  \left(  1-q^{2n}\right)  \right)  =\sum_{\ell\in\mathbb{Z}%
}q^{\ell^{2}}z^{\ell}.
\]

\end{theorem}

However, we aren't just planning to view this identity as a formal identity
between power series; instead, we will later evaluate both sides at certain
powers of another indeterminate $x$ (i.e., we will set $q=x^{a}$ and $z=x^{b}$
for some positive integers $a$ and $b$). Alas, this is not an operation
defined on the whole ring $\left(  \mathbb{Z}\left[  z^{\pm}\right]  \right)
\left[  \left[  q\right]  \right]  $. For example, setting $q=x$ and $z=x$ in
the sum $\sum_{\ell\in\mathbb{N}}q^{\ell}z^{-\ell}\in\left(  \mathbb{Z}\left[
z^{\pm}\right]  \right)  \left[  \left[  q\right]  \right]  $ yields the
nonsensical sum $\sum_{\ell\in\mathbb{N}}x^{\ell}x^{-\ell}=\sum_{\ell
\in\mathbb{Z}}1$.

Thus, Theorem \ref{thm.pars.jtp1} is not a version of Jacobi's triple product
identity that we can use for our purposes. Instead, let us interpret the
identity (\ref{eq.pars.jtp.jtp}) in a different way: Instead of treating $q$
and $z$ as indeterminates, I will set them to be powers of a single
indeterminate $x$ (more precisely, scalar multiples of such powers). This
leads us to the following version of the identity:

\begin{theorem}
[Jacobi's triple product identity, take 2]\label{thm.pars.jtp2}Let $a$ and $b$
be two integers such that $a>0$ and $a\geq\left\vert b\right\vert $. Let
$u,v\in\mathbb{Q}$ be rational numbers with $v\neq0$. In the ring
$\mathbb{Q}\left(  \left(  x\right)  \right)  $, set $q=ux^{a}$ and $z=vx^{b}%
$. Then,%
\[
\prod_{n>0}\left(  \left(  1+q^{2n-1}z\right)  \left(  1+q^{2n-1}%
z^{-1}\right)  \left(  1-q^{2n}\right)  \right)  =\sum_{\ell\in\mathbb{Z}%
}q^{\ell^{2}}z^{\ell}.
\]

\end{theorem}

Before we start proving this theorem, let us check that the infinite product
on its left hand side and the infinite sum on its right are well-defined:

\begin{itemize}
\item The infinite product is%
\begin{align*}
&  \prod_{n>0}\left(  \left(  1+q^{2n-1}z\right)  \left(  1+q^{2n-1}%
z^{-1}\right)  \left(  1-q^{2n}\right)  \right) \\
&  =\prod_{n>0}\left(  \left(  1+\left(  ux^{a}\right)  ^{2n-1}\left(
vx^{b}\right)  \right)  \left(  1+\left(  ux^{a}\right)  ^{2n-1}\left(
vx^{b}\right)  ^{-1}\right)  \left(  1-\left(  ux^{a}\right)  ^{2n}\right)
\right) \\
&  =\prod_{n>0}\left(  \left(  1+u^{2n-1}vx^{\left(  2n-1\right)  a+b}\right)
\left(  1+u^{2n-1}v^{-1}x^{\left(  2n-1\right)  a-b}\right)  \left(
1-u^{2n}x^{2na}\right)  \right)  .
\end{align*}
All factors in this product belong to the ring $\mathbb{Q}\left[  \left[
x\right]  \right]  $ (not just to $\mathbb{Q}\left(  \left(  x\right)
\right)  $), since the exponents $\left(  2n-1\right)  a+b$ and $\left(
2n-1\right)  a-b$ and $2na$ are always nonnegative for any $n>0$ (indeed, for
any $n>0$, we have $\underbrace{\left(  2n-1\right)  }_{\geq1}\underbrace{a}%
_{\geq\left\vert b\right\vert }+\,b\geq\left\vert b\right\vert +b\geq0$ and
$\underbrace{\left(  2n-1\right)  }_{\geq1}\underbrace{a}_{\geq\left\vert
b\right\vert }-\,b\geq\left\vert b\right\vert -b\geq0$ and $2n\underbrace{a}%
_{>0}>0$). Moreover, this product is multipliable, because

\begin{itemize}
\item the number $\left(  2n-1\right)  a+b$ grows linearly when $n\rightarrow
\infty$ (since $a>0$);

\item the number $\left(  2n-1\right)  a-b$ grows linearly when $n\rightarrow
\infty$ (since $a>0$);

\item the number $2na$ grows linearly when $n\rightarrow\infty$ (since $a>0$).
\end{itemize}

Thus, the infinite product is well-defined.

This argument also shows that the three subproducts $\prod_{n>0}\left(
1+q^{2n-1}z\right)  $ and $\prod_{n>0}\left(  1+q^{2n-1}z^{-1}\right)  $ and
$\prod_{n>0}\left(  1-q^{2n}\right)  $ are well-defined.

\item The infinite sum is%
\[
\sum_{\ell\in\mathbb{Z}}q^{\ell^{2}}z^{\ell}=\sum_{\ell\in\mathbb{Z}}\left(
ux^{a}\right)  ^{\ell^{2}}\left(  vx^{b}\right)  ^{\ell}=\sum_{\ell
\in\mathbb{Z}}u^{\ell^{2}}v^{\ell}x^{a\ell^{2}+b\ell}.
\]
All addends in this sum belong to the ring $\mathbb{Q}\left[  \left[
x\right]  \right]  $ (not just to $\mathbb{Q}\left(  \left(  x\right)
\right)  $), since the exponent $a\ell^{2}+b\ell$ is always nonnegative for
any $\ell\in\mathbb{Z}$ (indeed, we have $\underbrace{a}_{\geq\left\vert
b\right\vert }\underbrace{\ell^{2}}_{=\left\vert \ell\right\vert ^{2}%
}+\underbrace{b\ell}_{\geq-\left\vert b\ell\right\vert =-\left\vert
b\right\vert \cdot\left\vert \ell\right\vert }\geq\left\vert b\right\vert
\cdot\left\vert \ell\right\vert ^{2}-\left\vert b\right\vert \cdot\left\vert
\ell\right\vert =\underbrace{\left\vert b\right\vert }_{\geq0}\cdot
\underbrace{\left\vert \ell\right\vert \cdot\left(  \left\vert \ell\right\vert
-1\right)  }_{\geq0}\geq0$ for any $\ell\in\mathbb{Z}$). Moreover, this sum is
summable, because $a\ell^{2}+b\ell$ grows quadratically when $\ell
\rightarrow+\infty$ or $\ell\rightarrow-\infty$ (since $a>0$).
\end{itemize}

\subsubsection{Jacobi implies Euler}

We will give a proof of Jacobi's triple product identity that works equally
for both versions of it (Theorem \ref{thm.pars.jtp1} and Theorem
\ref{thm.pars.jtp2}). But first, let us see how it yields Euler's pentagonal
number theorem as a particular case.

\begin{proof}
[Proof of Theorem \ref{thm.pars.pent} using Theorem \ref{thm.pars.jtp2}.]Set
$q=x^{3}$ and $z=-x$ in Theorem \ref{thm.pars.jtp2}. (This means that we apply
Theorem \ref{thm.pars.jtp2} to $a=3$ and $b=1$ and $u=1$ and $v=-1$). We get%
\begin{align}
&  \prod_{n>0}\left(  \left(  1+\left(  x^{3}\right)  ^{2n-1}\left(
-x\right)  \right)  \left(  1+\left(  x^{3}\right)  ^{2n-1}\left(  -x\right)
^{-1}\right)  \left(  1-\left(  x^{3}\right)  ^{2n}\right)  \right)
\nonumber\\
&  =\sum_{\ell\in\mathbb{Z}}\left(  x^{3}\right)  ^{\ell^{2}}\left(
-x\right)  ^{\ell}. \label{pf.thm.pars.pent.1}%
\end{align}
The left hand side of this equality simplifies as follows:%
\begin{align*}
&  \prod_{n>0}\left(  \left(  1+\underbrace{\left(  x^{3}\right)
^{2n-1}\left(  -x\right)  }_{\substack{=-x^{3\left(  2n-1\right)
+1}\\=-x^{6n-2}}}\right)  \left(  1+\underbrace{\left(  x^{3}\right)
^{2n-1}\left(  -x\right)  ^{-1}}_{\substack{=-x^{3\left(  2n-1\right)
-1}\\=-x^{6n-4}}}\right)  \left(  1-\underbrace{\left(  x^{3}\right)  ^{2n}%
}_{=x^{6n}}\right)  \right) \\
&  =\prod_{n>0}\left(  \left(  1-\underbrace{x^{6n-2}}_{=\left(  x^{2}\right)
^{3n-1}}\right)  \left(  1-\underbrace{x^{6n-4}}_{=\left(  x^{2}\right)
^{3n-2}}\right)  \left(  1-\underbrace{x^{6n}}_{=\left(  x^{2}\right)  ^{3n}%
}\right)  \right) \\
&  =\prod_{n>0}\left(  \left(  1-\left(  x^{2}\right)  ^{3n-1}\right)  \left(
1-\left(  x^{2}\right)  ^{3n-2}\right)  \left(  1-\left(  x^{2}\right)
^{3n}\right)  \right) \\
&  =\prod_{k>0}\left(  1-\left(  x^{2}\right)  ^{k}\right)  ,
\end{align*}
since each positive integer $k$ can be uniquely represented as $3n-1$ or
$3n-2$ or $3n$ for some positive integer $n$.

Comparing this with (\ref{pf.thm.pars.pent.1}), we obtain%
\begin{align}
&  \prod_{k>0}\left(  1-\left(  x^{2}\right)  ^{k}\right) \nonumber\\
&  =\sum_{\ell\in\mathbb{Z}}\underbrace{\left(  x^{3}\right)  ^{\ell^{2}%
}\left(  -x\right)  ^{\ell}}_{=\left(  -1\right)  ^{\ell}x^{3\ell^{2}+\ell}%
}=\sum_{\ell\in\mathbb{Z}}\left(  -1\right)  ^{\ell}\underbrace{x^{3\ell
^{2}+\ell}}_{\substack{=x^{2w_{-\ell}}\\\text{(since }3\ell^{2}+\ell=\left(
3\ell+1\right)  \ell=\left(  3\left(  -\ell\right)  -1\right)  \left(
-\ell\right)  =2w_{-\ell}\\\text{(because }w_{-\ell}\text{ is defined as
}\left(  3\left(  -\ell\right)  -1\right)  \left(  -\ell\right)  /2\text{))}%
}}\nonumber\\
&  =\sum_{\ell\in\mathbb{Z}}\underbrace{\left(  -1\right)  ^{\ell}}_{=\left(
-1\right)  ^{-\ell}}\underbrace{x^{2w_{-\ell}}}_{=\left(  x^{2}\right)
^{w_{-\ell}}}=\sum_{\ell\in\mathbb{Z}}\left(  -1\right)  ^{-\ell}\left(
x^{2}\right)  ^{w_{-\ell}}\nonumber\\
&  =\sum_{k\in\mathbb{Z}}\left(  -1\right)  ^{k}\left(  x^{2}\right)  ^{w_{k}}
\label{pf.thm.pars.pent.2}%
\end{align}
(here, we have substituted $k$ for $-\ell$ in the sum).

Now, let us \textquotedblleft substitute $x$ for $x^{2}$\textquotedblright\ in
this equality (see below for how this works). As a result, we obtain%
\[
\prod_{k>0}\left(  1-x^{k}\right)  =\sum_{k\in\mathbb{Z}}\left(  -1\right)
^{k}x^{w_{k}}.
\]
This is Euler's pentagonal number theorem (Theorem \ref{thm.pars.pent}).
\end{proof}

What did I mean by \textquotedblleft substituting $x$ for $x^{2}%
$\textquotedblright? I meant using the following simple fact:

\begin{lemma}
\label{lem.fps.fxx=gxx}Let $K$ be a commutative ring. Let $f$ and $g$ be two
FPSs in $K\left[  \left[  x\right]  \right]  $. Assume that $f\left[
x^{2}\right]  =g\left[  x^{2}\right]  $. Then, $f=g$.
\end{lemma}

\begin{proof}
This is easy: Write $f$ and $g$ as $f=\sum_{n\in\mathbb{N}}f_{n}x^{n}$ and
$g=\sum_{n\in\mathbb{N}}g_{n}x^{n}$ where $f_{0},f_{1},f_{2},\ldots\in K$ and
$g_{0},g_{1},g_{2},\ldots\in K$. Then, $f\left[  x^{2}\right]  =\sum
_{n\in\mathbb{N}}f_{n}\left(  x^{2}\right)  ^{n}=\sum_{n\in\mathbb{N}}%
f_{n}x^{2n}$ and similarly $g\left[  x^{2}\right]  =\sum_{n\in\mathbb{N}}%
g_{n}x^{2n}$. Thus, our assumption $f\left[  x^{2}\right]  =g\left[
x^{2}\right]  $ rewrites as $\sum_{n\in\mathbb{N}}f_{n}x^{2n}=\sum
_{n\in\mathbb{N}}g_{n}x^{2n}$. Comparing $x^{2n}$-coefficients in this
equality, we conclude that $f_{n}=g_{n}$ for each $n\in\mathbb{N}$. Hence,
$\sum_{n\in\mathbb{N}}\underbrace{f_{n}}_{=g_{n}}x^{n}=\sum_{n\in\mathbb{N}%
}g_{n}x^{n}$. In other words, $f=g$ (since $f=\sum_{n\in\mathbb{N}}f_{n}x^{n}$
and $g=\sum_{n\in\mathbb{N}}g_{n}x^{n}$). This proves Lemma
\ref{lem.fps.fxx=gxx}.
\end{proof}

Lemma \ref{lem.fps.fxx=gxx} justifies our \textquotedblleft substituting $x$
for $x^{2}$\textquotedblright\ in the above proof; indeed, we can apply Lemma
\ref{lem.fps.fxx=gxx} to $K=\mathbb{Q}$ and $f=\prod_{k>0}\left(
1-x^{k}\right)  $ and $g=\sum_{k\in\mathbb{Z}}\left(  -1\right)  ^{k}x^{w_{k}%
}$ (because (\ref{pf.thm.pars.pent.2}) says that these two FPSs $f$ and $g$
satisfy $f\left[  x^{2}\right]  =g\left[  x^{2}\right]  $), and consequently
obtain $\prod_{k>0}\left(  1-x^{k}\right)  =\sum_{k\in\mathbb{Z}}\left(
-1\right)  ^{k}x^{w_{k}}$. Thus, Theorem \ref{thm.pars.pent} is proved using
Theorem \ref{thm.pars.jtp2}. It therefore remains to prove the latter.

\subsubsection{Proof of Jacobi's triple product identity}

The following proof is due to Borcherds, and I have taken it from
\cite[\S 8.3]{Camero16} (note that \cite[\S 11.2]{Loehr-BC} gives essentially
the same proof, albeit in a different language).

\begin{proof}
[Proof of Theorem \ref{thm.pars.jtp1} and Theorem \ref{thm.pars.jtp2}.]The
following argument applies equally to Theorem \ref{thm.pars.jtp1} and to
Theorem \ref{thm.pars.jtp2}. (The meanings of $q$ and $z$ differ between these
two theorems, but all the infinite sums and products considered below are
well-defined in either case.)

We will use a somewhat physics-inspired language:

\begin{itemize}
\item A \emph{level} will mean a number of the form $p+\dfrac{1}{2}$ with
$p\in\mathbb{Z}$. (Thus, there is exactly one level midway between any two
consecutive integers.)

\item A \emph{state} will mean a set of levels that contains

\begin{itemize}
\item all but finitely many negative levels, and

\item only finitely many positive levels.
\end{itemize}
\end{itemize}

Here is an example of a state:%
\[
\left\{  \dfrac{-5}{2},\ \ \dfrac{-1}{2},\ \ \dfrac{1}{2},\ \ \dfrac{3}%
{2},\ \ \dfrac{7}{2},\ \ \dfrac{13}{2}\right\}  \cup\left\{  \text{all levels
}\leq\dfrac{-9}{2}\right\}  .
\]
Visually, it can be represented as follows:%
\[%
% [tikz figure removed]
%BeginExpansion
% [tikz figure removed]%
%EndExpansion
\]
where

\begin{itemize}
\item A white (=hollow) circle
%TCIMACRO{\TeXButton{tikz hollow circle}{\raisebox{-0.5ex}{\tikz\draw
%[black,fill=white] (0,0) circle (1ex);}}}%
%BeginExpansion
\raisebox{-0.5ex}{\tikz\draw[black,fill=white] (0,0) circle (1ex);}%
%EndExpansion
\ means a level that is contained in the state (you can think of it as an
\textquotedblleft electron\textquotedblright).

\item A black (=filled) circle
%TCIMACRO{\TeXButton{tikz full circle}{\raisebox{-0.5ex}{\tikz\draw
%[black,fill=black] (0,0) circle (1ex);}} }%
%BeginExpansion
\raisebox{-0.5ex}{\tikz\draw[black,fill=black] (0,0) circle (1ex);}
%EndExpansion
means a level that is not contained in the state (think of it as a
\textquotedblleft hole\textquotedblright).
\end{itemize}

For any state $S$,

\begin{itemize}
\item we define the \emph{energy} of $S$ to be%
\[
\operatorname*{energy}S:=\sum_{\substack{p>0;\\p\in S}}\underbrace{2p}%
_{>0}-\sum_{\substack{p<0;\\p\notin S}}\underbrace{2p}_{<0}\in\mathbb{N}%
\]
(where the summation index $p$ in the first sum runs over the finitely many
positive levels contained in $S$, while the summation index $p$ in the second
sum runs over the finitely many negative levels not contained in $S$).

\item we define the \emph{particle number} of $S$ to be%
\begin{align*}
\operatorname*{parnum}S:=  &  \left(  \text{\# of levels }p>0\text{ such that
}p\in S\right) \\
&  \ \ \ \ \ \ \ \ \ \ -\left(  \text{\# of levels }p<0\text{ such that
}p\notin S\right)  \in\mathbb{Z}.
\end{align*}

\end{itemize}

For instance, in the above example, we have%
\[
\operatorname*{energy}S=1+3+7+13-\left(  -3\right)  -\left(  -7\right)  =34
\]
and%
\[
\operatorname*{parnum}S=4-2=2.
\]


We want to prove the identity%
\[
\prod_{n>0}\left(  \left(  1+q^{2n-1}z\right)  \left(  1+q^{2n-1}%
z^{-1}\right)  \left(  1-q^{2n}\right)  \right)  =\sum_{\ell\in\mathbb{Z}%
}q^{\ell^{2}}z^{\ell}.
\]
We will first transform this identity into an equivalent one: Namely, we move
the $\left(  1-q^{2n}\right)  $ factors from the left hand side to the right
hand side by multiplying both sides with $\prod_{n>0}\left(  1-q^{2n}\right)
^{-1}$. Thus, we can rewrite our identity as%
\[
\prod_{n>0}\left(  \left(  1+q^{2n-1}z\right)  \left(  1+q^{2n-1}%
z^{-1}\right)  \right)  =\left(  \sum_{\ell\in\mathbb{Z}}q^{\ell^{2}}z^{\ell
}\right)  \prod_{n>0}\left(  1-q^{2n}\right)  ^{-1}.
\]
We will prove this new identity by showing that both of its sides are%
\[
\sum_{S\text{ is a state}}q^{\operatorname*{energy}S}z^{\operatorname*{parnum}%
S}.
\]


\textit{Left hand side:} We have%
\begin{align}
&  \prod_{n>0}\left(  \left(  1+q^{2n-1}z\right)  \left(  1+q^{2n-1}%
z^{-1}\right)  \right) \nonumber\\
&  =\left(  \prod_{n>0}\left(  1+q^{2n-1}z\right)  \right)  \left(
\prod_{n>0}\left(  1+q^{2n-1}z^{-1}\right)  \right) \nonumber\\
&  =\left(  \prod_{p\text{ is a positive level}}\left(  1+q^{2p}z\right)
\right)  \left(  \prod_{p\text{ is a negative level}}\left(  1+q^{-2p}%
z^{-1}\right)  \right) \nonumber\\
&  \ \ \ \ \ \ \ \ \ \ \ \ \ \ \ \ \ \ \ \ \left(
\begin{array}
[c]{c}%
\text{here, we have substituted }p+\dfrac{1}{2}\text{ for }n\text{ in the
first product,}\\
\text{and have substituted }-p+\dfrac{1}{2}\text{ for }n\text{ in the second
product}%
\end{array}
\right) \nonumber\\
&  =\left(  \sum_{\substack{P\text{ is a finite set}\\\text{of positive
levels}}}\ \ \prod_{p\in P}\left(  q^{2p}z\right)  \right)  \left(
\sum_{\substack{N\text{ is a finite set}\\\text{of negative levels}}%
}\ \ \prod_{p\in N}\left(  q^{-2p}z^{-1}\right)  \right) \nonumber\\
&  \ \ \ \ \ \ \ \ \ \ \ \ \ \ \ \ \ \ \ \ \left(  \text{here, we have
expanded both products using (\ref{eq.fps.prod.binary.prod-inf2})}\right)
\nonumber\\
&  =\sum_{\substack{P\text{ is a finite set}\\\text{of positive levels}%
}}\ \ \sum_{\substack{N\text{ is a finite set}\\\text{of negative levels}%
}}\ \ \underbrace{\prod_{p\in P}\left(  q^{2p}z\right)  \ \ \prod_{p\in
N}\left(  q^{-2p}z^{-1}\right)  }_{=q^{2\left(  \text{sum of elements of
}P\right)  -2\left(  \text{sum of elements of }N\right)  }z^{\left\vert
P\right\vert -\left\vert N\right\vert }}\nonumber\\
&  =\sum_{\substack{P\text{ is a finite set}\\\text{of positive levels}%
}}\ \ \sum_{\substack{N\text{ is a finite set}\\\text{of negative levels}%
}}q^{2\left(  \text{sum of elements of }P\right)  -2\left(  \text{sum of
elements of }N\right)  }z^{\left\vert P\right\vert -\left\vert N\right\vert
}\nonumber\\
&  =\sum_{S\text{ is a state}}\underbrace{q^{2\left(  \text{sum of positive
levels in }S\right)  -2\left(  \text{sum of negative levels not in }S\right)
}}_{\substack{=q^{\operatorname*{energy}S}\\\text{(by the definition of
}\operatorname*{energy}S\text{)}}}\nonumber\\
&  \ \ \ \ \ \ \ \ \ \ \ \ \ \ \ \ \ \ \ \ \underbrace{z^{\left(  \text{number
of positive levels in }S\right)  -\left(  \text{number of negative levels not
in }S\right)  }}_{\substack{=z^{\operatorname*{parnum}S}\\\text{(by the
definition of }\operatorname*{parnum}S\text{)}}}\nonumber\\
&  \ \ \ \ \ \ \ \ \ \ \ \ \ \ \ \ \ \ \ \ \ \ \ \ \ \ \ \ \ \ \left(
\begin{array}
[c]{c}%
\text{here, we have combined }P\text{ and }N\\
\text{into a single state }S:=P\cup\overline{N}\text{,}\\
\text{where }\overline{N}=\left\{  \text{all negative levels}\right\}
\setminus N
\end{array}
\right) \nonumber\\
&  =\sum_{S\text{ is a state}}q^{\operatorname*{energy}S}%
z^{\operatorname*{parnum}S}. \label{pf.thm.pars.jtp2.lhs}%
\end{align}


\textit{Right hand side:} Recall that%
\begin{align*}
\prod_{n>0}\left(  1-x^{n}\right)  ^{-1}  &  =\prod_{n>0}\dfrac{1}{1-x^{n}%
}=\sum_{n\in\mathbb{N}}p\left(  n\right)  x^{n}\ \ \ \ \ \ \ \ \ \ \left(
\text{by Theorem \ref{thm.pars.main-gf}}\right) \\
&  =\sum_{\substack{\lambda\text{ is a}\\\text{partition}}}x^{\left\vert
\lambda\right\vert }%
\end{align*}
(because the sum $\sum_{\substack{\lambda\text{ is a}\\\text{partition}%
}}x^{\left\vert \lambda\right\vert }$ contains each monomial $x^{n}$ precisely
$p\left(  n\right)  $ times). Substituting $q^{2}$ for $x$ in this equality,
we find%
\[
\prod_{n>0}\left(  1-\left(  q^{2}\right)  ^{n}\right)  ^{-1}=\sum
_{\substack{\lambda\text{ is a}\\\text{partition}}}\left(  q^{2}\right)
^{\left\vert \lambda\right\vert }.
\]
In other words,%
\[
\prod_{n>0}\left(  1-q^{2n}\right)  ^{-1}=\sum_{\substack{\lambda\text{ is
a}\\\text{partition}}}q^{2\left\vert \lambda\right\vert }.
\]
Multiplying both sides of this equality by $\sum_{\ell\in\mathbb{Z}}%
q^{\ell^{2}}z^{\ell}$, we obtain%
\begin{align}
\left(  \sum_{\ell\in\mathbb{Z}}q^{\ell^{2}}z^{\ell}\right)  \prod
_{n>0}\left(  1-q^{2n}\right)  ^{-1}  &  =\left(  \sum_{\ell\in\mathbb{Z}%
}q^{\ell^{2}}z^{\ell}\right)  \sum_{\substack{\lambda\text{ is a}%
\\\text{partition}}}q^{2\left\vert \lambda\right\vert }\nonumber\\
&  =\sum_{\ell\in\mathbb{Z}}\ \ \sum_{\substack{\lambda\text{ is
a}\\\text{partition}}}q^{\ell^{2}+2\left\vert \lambda\right\vert }z^{\ell}.
\label{pf.thm.pars.jtp2.rhs1}%
\end{align}
We want to show that this equals $\sum\limits_{S\text{ is a state}%
}q^{\operatorname*{energy}S}z^{\operatorname*{parnum}S}$. In order to do this,
we will find a bijection%
\[
\Phi_{\ell}:\left\{  \text{partitions}\right\}  \rightarrow\left\{
\text{states with particle number }\ell\right\}
\]
for each $\ell\in\mathbb{Z}$, and we will show that this bijection satisfies
\[
\operatorname*{energy}\left(  \Phi_{\ell}\left(  \lambda\right)  \right)
=\ell^{2}+2\left\vert \lambda\right\vert \ \ \ \ \ \ \ \ \ \ \text{for each
}\ell\in\mathbb{Z}\text{ and }\lambda\in\left\{  \text{partitions}\right\}  .
\]


Let us do this. Fix $\ell\in\mathbb{Z}$. We define the state $G_{\ell}$
(called the \textquotedblleft$\ell$-ground state\textquotedblright) by%
\[
G_{\ell}:=\left\{  \text{all levels }<\ell\right\}  =\left\{  \ell-\dfrac
{1}{2},\ \ \ell-\dfrac{3}{2},\ \ \ell-\dfrac{5}{2},\ \ \ldots\right\}  .
\]
Here is how it looks like (for $\ell$ positive):%
\[%
% [tikz figure removed]
%BeginExpansion
% [tikz figure removed]%
%EndExpansion
\]


If $\ell\geq0$, then this state $G_{\ell}$ has energy
\[
\operatorname*{energy}G_{\ell}=1+3+5+\cdots+\left(  2\ell-1\right)  =\ell^{2}%
\]
and particle number%
\[
\operatorname*{parnum}G_{\ell}=\ell-0=\ell.
\]
If $\ell<0$, then it has energy%
\[
\operatorname*{energy}G_{\ell}=-\left(  -1\right)  -\left(  -3\right)
-\left(  -5\right)  -\cdots-\left(  2\ell+1\right)  =\ell^{2}%
\]
and particle number%
\[
\operatorname*{parnum}G_{\ell}=0-\left(  -\ell\right)  =\ell.
\]
Note that the answers are the same in both cases. Thus, whatever sign $\ell$
has, we have%
\[
\operatorname*{energy}G_{\ell}=\ell^{2}\ \ \ \ \ \ \ \ \ \ \text{and}%
\ \ \ \ \ \ \ \ \ \ \operatorname*{parnum}G_{\ell}=\ell.
\]


If $S$ is a state, and if $p\in S$, and if $q$ is a positive integer such that
$p+q\notin S$, then we define a new state%
\[
\operatorname*{jump}\nolimits_{p,q}S:=\left(  S\setminus\left\{  p\right\}
\right)  \cup\left\{  p+q\right\}  .
\]
We say that this state $\operatorname*{jump}\nolimits_{p,q}S$ is obtained from
$S$ by letting the electron at level $p$ \emph{jump} $q$ steps to the right.
Note that $\operatorname*{jump}\nolimits_{p,q}S$ has the same particle number
as $S$ (check this!\footnote{There are three cases:
\par
\textit{Case 1:} We have $p>0$ (that is, the particle jumps from a positive
level to a positive level).
\par
\textit{Case 2:} We have $p<0$ and $p+q>0$ (that is, the particle jumps from a
negative level to a positive level).
\par
\textit{Case 3:} We have $p+q<0$ (that is, the particle jumps from a negative
level to a negative level).
\par
Each case is easy to check.}), whereas its energy is $2q$ higher than that of
$S$ (check this!\footnote{Again, the same three cases are to be considered.}).
Thus, a jumping electron raises the energy but keeps the particle number unchanged.

For any partition $\lambda=\left(  \lambda_{1},\lambda_{2},\ldots,\lambda
_{k}\right)  $, we define the state $E_{\ell,\lambda}$ (called an
\textquotedblleft excited state\textquotedblright) by starting with the $\ell
$-ground state $G_{\ell}$, and then successively letting the $k$ electrons at
the highest levels (which are -- from highest to lowest -- the levels
$\ell-1+\dfrac{1}{2},\ \ \ell-2+\dfrac{1}{2},\ \ \ldots,\ \ \ell-k+\dfrac
{1}{2}$) jump $\lambda_{1},\lambda_{2},\ldots,\lambda_{k}$ steps to the right,
respectively (starting with the rightmost electron). In other words,%
\begin{align*}
E_{\ell,\lambda}  &  :=\operatorname*{jump}\nolimits_{\ell-k+1/2,\ \ \lambda
_{k}}\left(  \cdots\left(  \operatorname*{jump}\nolimits_{\ell
-2+1/2,\ \ \lambda_{2}}\left(  \operatorname*{jump}\nolimits_{\ell
-1+1/2,\ \ \lambda_{1}}\left(  G_{\ell}\right)  \right)  \right)  \right) \\
&  =\left\{  \text{all levels }<\ell-k\right\}  \cup\left\{  \ell-i+\dfrac
{1}{2}+\lambda_{i}\ \mid\ i\in\left\{  1,2,\ldots,k\right\}  \right\}  .
\end{align*}
(Check that these jumps are well-defined -- i.e., that each electron jumps to
an unoccupied state.)

[Example: Let $\ell=3$ and $k=4$ and $\lambda=\left(  \lambda_{1},\lambda
_{2},\lambda_{3},\lambda_{4}\right)  =\left(  4,2,2,1\right)  $. Then,%
\[
E_{\ell,\lambda}=\left\{  \text{all levels }<-1\right\}  \cup\left\{
\dfrac{3}{2},\ \ \dfrac{7}{2},\ \ \dfrac{9}{2},\ \ \dfrac{13}{2}\right\}  .
\]
Here is a picture of this state $E_{\ell,\lambda}$ and how it is constructed
by a sequence of electron jumps (the top state is $G_{\ell}$; the bottom state
is $E_{\ell,\lambda}$):%
\[%
% [tikz figure removed]
%BeginExpansion
% [tikz figure removed]%
%EndExpansion
.
\]
Note how the order of the electrons after the jumps is the same as before --
i.e., the electron in the rightmost position before the jumps is still the
rightmost one after the jumps, etc.]

Recall that the $\ell$-ground state $G_{\ell}$ has energy $\ell^{2}$ and
particle number $\ell$. The state $E_{\ell,\lambda}$ that we have just defined
is obtained from this $\ell$-ground state $G_{\ell}$ by $k$ jumps, which are
jumps by $\lambda_{1},\lambda_{2},\ldots,\lambda_{k}$ steps respectively.
Recall that a jump by $q$ steps raises the energy by $2q$ but keeps the
particle number unchanged. Thus, the \textquotedblleft excited
state\textquotedblright\ $E_{\ell,\lambda}$ has energy $\ell^{2}%
+\underbrace{2\lambda_{1}+2\lambda_{2}+\cdots+2\lambda_{k}}_{=2\left\vert
\lambda\right\vert }=\ell^{2}+2\left\vert \lambda\right\vert $ and particle
number $\ell$. Furthermore, every state with particle number $\ell$ can be
written as $E_{\ell,\lambda}$ for a unique partition $\lambda$ (check
this!\footnote{Here is how to obtain this $\lambda$: For each $i\in\left\{
1,2,3,\ldots\right\}  $, let $u_{i}$ be the $i$-th largest level in the state,
and let $\lambda_{i}=u_{i}-\ell+i-\dfrac{1}{2}$. Then, $\left(  \lambda
_{1},\lambda_{2},\lambda_{3},\ldots\right)  $ is a weakly decreasing sequence
of nonnegative integers, and all but finitely many of its entries are $0$
(since the state has particle number $\ell$, so that it is not hard to see
that $u_{i}=\ell-i+\dfrac{1}{2}$ for any sufficiently large $i$). Removing all
$0$s from this sequence $\left(  \lambda_{1},\lambda_{2},\lambda_{3}%
,\ldots\right)  $ thus results in a finite tuple $\left(  \lambda_{1}%
,\lambda_{2},\ldots,\lambda_{k}\right)  $, which is precisely the partition
$\lambda$ whose corresponding $E_{\ell,\lambda}$ is our state.}).

Thus, we obtain a bijection%
\begin{align*}
\Phi_{\ell}:\left\{  \text{partitions}\right\}   &  \rightarrow\left\{
\text{states with particle number }\ell\right\}  ,\\
\lambda &  \mapsto E_{\ell,\lambda}.
\end{align*}
This bijection satisfies%
\[
\operatorname*{energy}\left(  \Phi_{\ell}\left(  \lambda\right)  \right)
=\operatorname*{energy}E_{\ell,\lambda}=\ell^{2}+2\left\vert \lambda
\right\vert \ \ \ \ \ \ \ \ \ \ \text{for every partition }\lambda.
\]
Hence,%
\begin{align}
\sum_{\substack{\lambda\text{ is a}\\\text{partition}}}\underbrace{q^{\ell
^{2}+2\left\vert \lambda\right\vert }}_{\substack{=q^{\operatorname*{energy}%
\left(  \Phi_{\ell}\left(  \lambda\right)  \right)  }\\\text{(since
}\operatorname*{energy}\left(  \Phi_{\ell}\left(  \lambda\right)  \right)
=\ell^{2}+2\left\vert \lambda\right\vert \text{)}}}  &  =\sum
_{\substack{\lambda\text{ is a}\\\text{partition}}}q^{\operatorname*{energy}%
\left(  \Phi_{\ell}\left(  \lambda\right)  \right)  }\nonumber\\
&  =\sum_{\substack{S\text{ is a state with}\\\text{particle number }\ell
}}q^{\operatorname*{energy}S} \label{pf.thm.pars.jtp2.rhs5}%
\end{align}
(here, we have substituted $S$ for $\Phi_{\ell}\left(  \lambda\right)  $ in
the sum, since the map $\Phi_{\ell}$ is a bijection).

Forget that we fixed $\ell$. We thus have proved (\ref{pf.thm.pars.jtp2.rhs5})
for each $\ell\in\mathbb{Z}$.

Now, (\ref{pf.thm.pars.jtp2.rhs1}) becomes%
\begin{align*}
\left(  \sum_{\ell\in\mathbb{Z}}q^{\ell^{2}}z^{\ell}\right)  \prod
_{n>0}\left(  1-q^{2n}\right)  ^{-1}  &  =\sum_{\ell\in\mathbb{Z}%
}\ \ \underbrace{\sum_{\substack{\lambda\text{ is a}\\\text{partition}%
}}q^{\ell^{2}+2\left\vert \lambda\right\vert }}_{\substack{=\sum
_{\substack{S\text{ is a state with}\\\text{particle number }\ell
}}q^{\operatorname*{energy}S}\\\text{(by (\ref{pf.thm.pars.jtp2.rhs5}))}%
}}z^{\ell}\\
&  =\sum_{\ell\in\mathbb{Z}}\ \ \sum_{\substack{S\text{ is a state
with}\\\text{particle number }\ell}}q^{\operatorname*{energy}S}%
\underbrace{z^{\ell}}_{\substack{=z^{\operatorname*{parnum}S}\\\text{(since
}\ell=\operatorname*{parnum}S\text{)}}}\\
&  =\underbrace{\sum_{\ell\in\mathbb{Z}}\ \ \sum_{\substack{S\text{ is a state
with}\\\text{particle number }\ell}}}_{=\sum_{S\text{ is a state}}%
}q^{\operatorname*{energy}S}z^{\operatorname*{parnum}S}\\
&  =\sum_{S\text{ is a state}}q^{\operatorname*{energy}S}%
z^{\operatorname*{parnum}S}.
\end{align*}


Comparing this with (\ref{pf.thm.pars.jtp2.lhs}), we obtain%
\[
\prod_{n>0}\left(  \left(  1+q^{2n-1}z\right)  \left(  1+q^{2n-1}%
z^{-1}\right)  \right)  =\left(  \sum_{\ell\in\mathbb{Z}}q^{\ell^{2}}z^{\ell
}\right)  \prod_{n>0}\left(  1-q^{2n}\right)  ^{-1}.
\]
Multiplying both sides of this identity with $\prod_{n>0}\left(
1-q^{2n}\right)  $, we find%
\[
\prod_{n>0}\left(  \left(  1+q^{2n-1}z\right)  \left(  1+q^{2n-1}%
z^{-1}\right)  \left(  1-q^{2n}\right)  \right)  =\sum_{\ell\in\mathbb{Z}%
}q^{\ell^{2}}z^{\ell}.
\]
This proves Jacobi's Triple Product Identity (Theorem \ref{thm.pars.jtp1} and
Theorem \ref{thm.pars.jtp2}).
\end{proof}

Other proofs of Jacobi's Triple Product Identity can be found in
\cite[\S 3.4]{Aigner07}, \cite[Theorem 10.2]{Wagner08} and \cite[\S 1.3 and
\S 1.4]{Hirsch17}.

We note that proving Theorem \ref{thm.pars.pent} using Theorem
\ref{thm.pars.jtp2} can be viewed as a kind of overkill; there are more direct
proofs of Theorem \ref{thm.pars.pent} as well (see, e.g., \cite{Zabroc03} or
\cite[\S 10]{Koch16} or \cite[Exercise 5]{18f-mt3s} or \cite[\S 3]{Bell06}).

\subsubsection{Application: A recursion for the sum of divisors}

Now that Theorem \ref{thm.pars.pent} is proven, we use the occasion to
spotlight a truly unexpected application of it: Euler's recursion for the sum
of divisors. We state it in the following form:

\begin{theorem}
\label{thm.pars.euler-sum-div-rec}For any positive integer $n$, let
$\sigma\left(  n\right)  $ denote the sum of all positive divisors of $n$.
(For example, $\sigma\left(  6\right)  =1+2+3+6=12$ and $\sigma\left(
7\right)  =1+7=8$.)

Then, for each positive integer $n$, we have%
\[
\sum_{\substack{k\in\mathbb{Z};\\w_{k}<n}}\left(  -1\right)  ^{k}\sigma\left(
n-w_{k}\right)  =%
\begin{cases}
\left(  -1\right)  ^{k-1}n, & \text{if }n=w_{k}\text{ for some }k\in
\mathbb{Z};\\
0, & \text{if not.}%
\end{cases}
\]
(Here, $w_{k}$ is the $k$-th pentagonal number, defined in Definition
\ref{def.pars.pent-num}.)
\end{theorem}

The left hand side of the equality in Theorem \ref{thm.pars.euler-sum-div-rec}
looks as follows:%
\[
\sigma\left(  n\right)  -\sigma\left(  n-1\right)  -\sigma\left(  n-2\right)
+\sigma\left(  n-5\right)  +\sigma\left(  n-7\right)  -\sigma\left(
n-12\right)  -\sigma\left(  n-15\right)  \pm\cdots
\]
(the sum goes only as far as the arguments remain positive integers, so it is
a finite sum). Thus, by solving this equality for $\sigma\left(  n\right)  $,
we obtain a recursive formula for $\sigma\left(  n\right)  $. (In this form,
Theorem \ref{thm.pars.euler-sum-div-rec} appears in many places, such as
\cite[\S 23]{Ness61} and \cite[Theorem 37]{Johnso20}.) This does not give an
efficient algorithm for computing $\sigma\left(  n\right)  $ (it is much
slower than the formula in \cite[Exercise 2.18.1 \textbf{(b)}]{19s}, even
though the latter requires computing the prime factorization of $n$), but its
beauty and element of surprise (why would you expect the pentagonal numbers to
have anything to do with sums of divisors?) make up for this practical
uselessness. Even more surprisingly, it can be proved using Euler's pentagonal
number theorem, even though it has seemingly nothing to do with partitions!

\begin{proof}
[Proof of Theorem \ref{thm.pars.euler-sum-div-rec} (sketched).]First, let us
show that the right hand side in Theorem \ref{thm.pars.euler-sum-div-rec} is
well-defined. Indeed, it is easy to see that
\[
w_{0}<w_{1}<w_{-1}<w_{2}<w_{-2}<w_{3}<w_{-3}<\cdots;
\]
thus, the pentagonal numbers $w_{k}$ for all $k\in\mathbb{Z}$ are distinct.
Hence, if a positive integer $n$ satisfies $n=w_{k}$ for some $k\in\mathbb{Z}%
$, then this latter $k$ is uniquely determined by $n$, and therefore the
expression%
\[%
\begin{cases}
\left(  -1\right)  ^{k-1}n, & \text{if }n=w_{k}\text{ for some }k\in
\mathbb{Z};\\
0, & \text{if not}%
\end{cases}
\]
is well-defined. In other words, the right hand side in Theorem
\ref{thm.pars.euler-sum-div-rec} is well-defined.

Now, define the two FPSs $P$ and $S$ as in our above proof of Theorem
\ref{thm.pars.sigma1}. In that very proof, we have shown that these FPSs
satisfy
\begin{equation}
xP^{\prime}=SP. \label{pf.thm.pars.euler-sum-div-rec.xP'}%
\end{equation}


Next, we define a further FPS%
\[
Q:=\sum_{k\in\mathbb{Z}}\left(  -1\right)  ^{k}x^{w_{k}}\in\mathbb{Z}\left[
\left[  x\right]  \right]  .
\]


Multiplying the equalities%
\[
P=\sum_{n\in\mathbb{N}}p\left(  n\right)  x^{n}=\prod_{k=1}^{\infty}\dfrac
{1}{1-x^{k}}\ \ \ \ \ \ \ \ \ \ \left(  \text{by Theorem
\ref{thm.pars.main-gf}}\right)
\]
and%
\[
Q=\sum_{k\in\mathbb{Z}}\left(  -1\right)  ^{k}x^{w_{k}}=\prod\limits_{k=1}%
^{\infty}\left(  1-x^{k}\right)  \ \ \ \ \ \ \ \ \ \ \left(  \text{by Theorem
\ref{thm.pars.pent}}\right)  ,
\]
we obtain%
\[
PQ=\left(  \prod_{k=1}^{\infty}\dfrac{1}{1-x^{k}}\right)  \left(
\prod\limits_{k=1}^{\infty}\left(  1-x^{k}\right)  \right)  =\prod
\limits_{k=1}^{\infty}\underbrace{\left(  \dfrac{1}{1-x^{k}}\cdot\left(
1-x^{k}\right)  \right)  }_{=1}=1.
\]


Combining this equality with (\ref{pf.thm.pars.euler-sum-div-rec.xP'}), we can
easily see that $xQ^{\prime}=-QS$. Indeed, here are two (essentially
equivalent) ways to show this:

\begin{itemize}
\item From $PQ=1$, we obtain $Q=P^{-1}$. But $P\in\mathbb{Z}\left[  \left[
x\right]  \right]  _{1}$ (since the constant coefficient of $P$ is $p\left(
0\right)  =1$). Hence, Corollary \ref{cor.fps.loder.inv} (applied to $f=P$)
yields $\operatorname*{loder}\left(  P^{-1}\right)  =-\operatorname*{loder}P$.
In view of $P^{-1}=Q$, this rewrites as $\operatorname*{loder}%
Q=-\operatorname*{loder}P$. But the definition of a logarithmic derivative
yields $\operatorname*{loder}Q=\dfrac{Q^{\prime}}{Q}$ and
$\operatorname*{loder}P=\dfrac{P^{\prime}}{P}$. Hence,
\[
\dfrac{Q^{\prime}}{Q}=\operatorname*{loder}Q=-\operatorname*{loder}%
P=-\dfrac{P^{\prime}}{P}\ \ \ \ \ \ \ \ \ \ \left(  \text{since }%
\operatorname*{loder}P=\dfrac{P^{\prime}}{P}\right)  .
\]
Multiplying this equality by $x$, we find%
\begin{align*}
\dfrac{xQ^{\prime}}{Q}  &  =-\dfrac{xP^{\prime}}{P}=-\dfrac{SP}{P}%
\ \ \ \ \ \ \ \ \ \ \left(  \text{by (\ref{pf.thm.pars.euler-sum-div-rec.xP'}%
)}\right) \\
&  =-S,
\end{align*}
so that $xQ^{\prime}=-QS$.

\item A less conceptual but slicker way to draw the same conclusion is the
following: Taking derivatives in the equality $PQ=1$, we find $\left(
PQ\right)  ^{\prime}=1^{\prime}=0$, so that $0=\left(  PQ\right)  ^{\prime
}=P^{\prime}Q+PQ^{\prime}$ (by Theorem \ref{thm.fps.deriv.rules}
\textbf{(d)}). Thus, $PQ^{\prime}=-P^{\prime}Q$. Multiplying this equality by
$x$, we find $PQ^{\prime}\cdot x=-P^{\prime}Q\cdot x=-\underbrace{xP^{\prime}%
}_{\substack{=SP\\\text{(by (\ref{pf.thm.pars.euler-sum-div-rec.xP'}))}%
}}Q=-S\underbrace{PQ}_{=1}=-S$. Multiplying this further by $Q$, we obtain
$PQ^{\prime}\cdot xQ=-SQ=-QS$. In view of $PQ^{\prime}\cdot xQ=xQ^{\prime
}\cdot\underbrace{PQ}_{=1}=xQ^{\prime}$, we can rewrite this as $xQ^{\prime
}=-QS$.
\end{itemize}

Either way, we now know that $xQ^{\prime}=-QS$. However, from $Q=\sum
_{k\in\mathbb{Z}}\left(  -1\right)  ^{k}x^{w_{k}}$, we obtain%
\[
Q^{\prime}=\left(  \sum_{k\in\mathbb{Z}}\left(  -1\right)  ^{k}x^{w_{k}%
}\right)  ^{\prime}=\sum_{k\in\mathbb{Z}}\left(  -1\right)  ^{k}%
\underbrace{\left(  x^{w_{k}}\right)  ^{\prime}}_{\substack{=w_{k}x^{w_{k}%
-1}\\\text{(where we understand this}\\\text{expression to mean }0\text{ if
}w_{k}=0\text{)}}}=\sum_{k\in\mathbb{Z}}\left(  -1\right)  ^{k}w_{k}%
x^{w_{k}-1}.
\]
Upon multiplication by $x$, this becomes%
\[
xQ^{\prime}=x\sum_{k\in\mathbb{Z}}\left(  -1\right)  ^{k}w_{k}x^{w_{k}-1}%
=\sum_{k\in\mathbb{Z}}\left(  -1\right)  ^{k}w_{k}\underbrace{xx^{w_{k}-1}%
}_{=x^{w_{k}}}=\sum_{k\in\mathbb{Z}}\left(  -1\right)  ^{k}w_{k}x^{w_{k}}.
\]
Hence,%
\begin{align}
\sum_{k\in\mathbb{Z}}\left(  -1\right)  ^{k}w_{k}x^{w_{k}}  &  =xQ^{\prime
}=-\underbrace{Q}_{=\sum_{k\in\mathbb{Z}}\left(  -1\right)  ^{k}x^{w_{k}}%
}\ \ \underbrace{S}_{\substack{=\sum_{k>0}\sigma\left(  k\right)  x^{k}%
\\=\sum_{i>0}\sigma\left(  i\right)  x^{i}}}\nonumber\\
&  =-\left(  \sum_{k\in\mathbb{Z}}\left(  -1\right)  ^{k}x^{w_{k}}\right)
\left(  \sum_{i>0}\sigma\left(  i\right)  x^{i}\right) \nonumber\\
&  =\sum_{k\in\mathbb{Z}}\ \ \sum_{i>0}\underbrace{\left(  -\left(  -1\right)
^{k}\right)  }_{=\left(  -1\right)  ^{k-1}}\underbrace{x^{w_{k}}\sigma\left(
i\right)  x^{i}}_{=\sigma\left(  i\right)  x^{w_{k}+i}}\nonumber\\
&  =\sum_{k\in\mathbb{Z}}\ \ \sum_{i>0}\left(  -1\right)  ^{k-1}\sigma\left(
i\right)  x^{w_{k}+i}\nonumber\\
&  =\underbrace{\sum_{k\in\mathbb{Z}}\ \ \sum_{m>w_{k}}}_{=\sum_{m>0}%
\ \ \sum_{\substack{k\in\mathbb{Z};\\m>w_{k}}}}\left(  -1\right)  ^{k-1}%
\sigma\left(  m-w_{k}\right)  \underbrace{x^{w_{k}+\left(  m-w_{k}\right)  }%
}_{=x^{m}}\nonumber\\
&  \ \ \ \ \ \ \ \ \ \ \ \ \ \ \ \ \ \ \ \ \left(
\begin{array}
[c]{c}%
\text{here, we have substituted }m-w_{k}\text{ for }i\\
\text{in the inner sum}%
\end{array}
\right) \nonumber\\
&  =\sum_{m>0}\ \ \sum_{\substack{k\in\mathbb{Z};\\m>w_{k}}}\left(  -1\right)
^{k-1}\sigma\left(  m-w_{k}\right)  x^{m}.
\label{pf.thm.pars.euler-sum-div-rec.4}%
\end{align}


Now, let $n$ be a positive integer. Let us compare the coefficients of $x^{n}$
on the left and right hand sides of (\ref{pf.thm.pars.euler-sum-div-rec.4}).
On the left hand side, the coefficient of $x^{n}$ is%
\[%
\begin{cases}
\left(  -1\right)  ^{k}w_{k}, & \text{if }n=w_{k}\text{ for some }%
k\in\mathbb{Z};\\
0, & \text{if not}%
\end{cases}
\]
(since the pentagonal numbers $w_{k}$ for all $k\in\mathbb{Z}$ are distinct,
and thus the different addends on the left hand side of
(\ref{pf.thm.pars.euler-sum-div-rec.4}) contribute to different monomials). On
the right hand side, the coefficient of $x^{n}$ is obviously%
\[
\sum_{\substack{k\in\mathbb{Z};\\n>w_{k}}}\left(  -1\right)  ^{k-1}%
\sigma\left(  n-w_{k}\right)  .
\]
Since these two coefficients are equal (because
(\ref{pf.thm.pars.euler-sum-div-rec.4}) is an identity), we thus conclude that%
\begin{align*}%
\begin{cases}
\left(  -1\right)  ^{k}w_{k}, & \text{if }n=w_{k}\text{ for some }%
k\in\mathbb{Z};\\
0, & \text{if not}%
\end{cases}
&  =\underbrace{\sum_{\substack{k\in\mathbb{Z};\\n>w_{k}}}}_{=\sum
_{\substack{k\in\mathbb{Z};\\w_{k}<n}}}\underbrace{\left(  -1\right)  ^{k-1}%
}_{=-\left(  -1\right)  ^{k}}\sigma\left(  n-w_{k}\right) \\
&  =-\sum_{\substack{k\in\mathbb{Z};\\w_{k}<n}}\left(  -1\right)  ^{k}%
\sigma\left(  n-w_{k}\right)  .
\end{align*}
Thus,%
\begin{align*}
\sum_{\substack{k\in\mathbb{Z};\\w_{k}<n}}\left(  -1\right)  ^{k}\sigma\left(
n-w_{k}\right)   &  =-%
\begin{cases}
\left(  -1\right)  ^{k}w_{k}, & \text{if }n=w_{k}\text{ for some }%
k\in\mathbb{Z};\\
0, & \text{if not}%
\end{cases}
\\
&  =%
\begin{cases}
-\left(  -1\right)  ^{k}w_{k}, & \text{if }n=w_{k}\text{ for some }%
k\in\mathbb{Z};\\
0, & \text{if not}%
\end{cases}
\\
&  =%
\begin{cases}
\left(  -1\right)  ^{k-1}w_{k}, & \text{if }n=w_{k}\text{ for some }%
k\in\mathbb{Z};\\
0, & \text{if not}%
\end{cases}
\\
&  \ \ \ \ \ \ \ \ \ \ \ \ \ \ \ \ \ \ \ \ \left(  \text{since }-\left(
-1\right)  ^{k}=\left(  -1\right)  ^{k-1}\right) \\
&  =%
\begin{cases}
\left(  -1\right)  ^{k-1}n, & \text{if }n=w_{k}\text{ for some }k\in
\mathbb{Z};\\
0, & \text{if not}%
\end{cases}
\end{align*}
(because if $n=w_{k}$ for some $k\in\mathbb{Z}$, then $w_{k}=n$). This proves
Theorem \ref{thm.pars.euler-sum-div-rec}.
\end{proof}
